\chapter{分子动理论的应用}

\section{蒸发}

本章我们将进一步讨论分子动理论的一些应用。在前一章中，我们强调动理论的一个独特的方面，即分子或其他物体的任何一个自由度的平均动能都是$kT/2$。而我们现在将要讨论的中心则是，在不同位置处的单位体积中找到粒子的概率随$exp[-(\text{P.E.})/(kT)]$而变化。我们将讨论这个结论的一些应用。

我们想要研究的现象是十分复杂的：液体进行蒸发，金属中的电子逸出表面，或者包括大量原子的化学反应。由于情况过于复杂，对于这些情况不可能从分子动理论出发作出任何简单和正确的陈述。因此本章中除去特别强调的地方以外，都相应地不精确。要强调的概念只是：从分子动理论出发，我们能够或多或少理解事物的行为应当如何。利用热力学的论据，或由某些临界量的经验测量结果，我们就可以更准确地描写那些现象。

然而，即使只是或多或少知道一点为什么它的行为会那样，也是十分有用的，这样，当我们遇到某种新的情况，或者还没有着手分析过的情况时，我们也可以多少说出一点会发生什么事情。所以，尽管这里的讨论十分不精确，但本质上还是正确的，即在概念上是正确的，不过，比如说，在某些特殊细节上作了一点简化。

我们将考虑的第一个例子是液体的蒸发。假设有一个体积很大的容器，里面放有一部分在一定温度下和它的蒸气达到平衡状态的液体。假定蒸气的分子彼此相距很远，而在液体内的分子则紧靠在一起。现在的问题是要知道，与液体中的分子数相比较，气相的分子数有多少？在给定的温度下蒸气的密度有多大？它对温度的依赖关系如何？

我们不妨假设蒸气中的每单位体积的分子数为$n$。当然，分子数会随着温度而变化。假如对液体加热，蒸发增加，就得到更多的蒸气。现在用另一个量$1/V_a$表示液体中每单位体积的分子数。假定液体中的每个分子都占有一定的体积，若液体的分子较多，则合在一起所占的体积就较大。所以，若$V_a$为单个分子所占有的体积，那么单位体积除以每个分子的体积就是单位体积中的分子数。此外，我们假设，在分子之间存在着吸引力，以使它们能在液体中结合在一起，否则，我们就无法理解为什么液体是凝聚的。因而我们假设存在有这样一种力，并且在液体的分子中存在着束缚能，而当它变为蒸气时，就失去了这部分能量。这就是说，我们假定为了使一个分子从液体跑到蒸气中，必须对它作一定量的功W。一个处在蒸气状态下的分子和一个处在液体状态下的分子之间有一定的能量差W，因为我们是从其他分子对它的吸引中把它拉出来的。

现在我们利用上面的一般原理：在两个不同的区域中每单位体积的分子数的关系为
\begin{equation*}
    \frac{n_2}{n_1} = e^{-(E_2 - E_1)/kT}
\end{equation*}
因而蒸气中每单位体积的分子数$n$除以液体中每单位体积的分子数$1/V_a$，就等于
\begin{equation}
    \label{Eq:I:42:1}
    nV_a=e^{-W/kT}
\end{equation}
这是一条普遍的规则。它就像在重力作用下处于平衡状态的大气层，其下层气体要比上层气体稠密，因为把气体分子提到高度$h$ 处需要作$mgh$的功。液体中的分子比蒸气中的分子稠密，因为要通过能量的“小丘”$W$ 才能拉开它们，而两种密度比是$e^{-W/kT}$。

这就是我们所要推导的结果——蒸气密度随着$e$的某种负的能量除以$kT$的幂而变化。前面的因子我们实际上并不感兴趣，因为在大多数情况下蒸气密度远比液体密度小。在某些情况下，只要不是靠近临界点（那里两种密度几乎相同），而是蒸气密度远小于液体密度，$n$远小于$1/V_a$的事实是由$W$远大于$kT$的事实引起的。所以像式（\ref{Eq:I:42:1}）那样的公式只有在$W$大大超过$kT$的情况下才有意义，因为在这种情况下，我们把$e$自乘到一个很大的负指数幂，只要使$T$变化一点点，这个指数幂也变化一点点，而由指数因子引起的变化要比前面的因子可能出现的任何变化重要得多。为什么在像$V_a$这样的因子上会出现一些变化呢？因为我们的分析是近似的。说到底，实际上并不是每个分子都有一定的体积，当我们改变温度时，体积$V_a$不会保持不变——液体膨胀了。此外，还有与此类似的其他次要的特征，所以实际状况更为复杂。在整个过程中都有着随温度缓慢变化的因素。事实上可以说$W$本身也随温度而略有变化，因为在较高的温度和不同的分子体积下，有不同的平均吸引力等。因而，可以想象，如果我们有一个公式，其中每一个量都以一种未知的方式随温度而变化，那么实际上就等于完全没有公式。如果我们认识到，指数$W/kT$一般说来非常大，我们就会看出，作为温度函数的蒸气密度曲线，其主要的变化由指数因子引起，若取$W$为常数，以及取系数$1/V_a$近似为常数，对于曲线中比较短的间隔来说，这是一个很好的近似。换句话说，最主要的变化具有$e^{-W/kT}$这样的普遍特征。

结果表明，在自然界的许许多多种现象中都有这样一个特征：要从某个地方汲取一些能量，而它随温度变化的最主要的特征是有一个$e$的负能量除以$kT$的指数幂因子。只有当能量远大于$kT$，从而最主要的变化是由$kT$的变化引起，而不是由常数或其他因子引起时，这个特征才是一个有用的事实。

现在，我们考虑另一个会对蒸发得到大约相似结果的方法，但要对过程观察得更细致一些。在得出式（\ref{Eq:I:42:1}）时，我们曾简单地使用了在平衡状态下成立的一条规则，但为了更好地理解这一切，不妨观察一下所发生的具体细节。我们也可以以下述方式来描述所发生的事：蒸气中的分子不断地碰撞液体的表面，当它们撞击在液面上时，可能被弹开，也可能进入液体内。这里有一个未知的比例——可能是50对50，也可能是10对90——我们并不知道。不妨先认为它们都跑到液体内，以后再回过头来分析假设它们不全都跑到液体内的那种情况。于是在某一给定的时刻，总有一定数目的分子凝聚在液体表面。凝聚分子数，即到达单位面积上的分子数，等于每单位体积的数目$n$乘以速度$v$。分子的速度与温度有关，因为我们知道$mv^2/2$的平均值是$3kT/2$。所以$v$就是某种平均速度。当然，我们应当对各种方向的速度进行积分从而得到某种平均值，但粗略地说它正比于方均根速度，并带有某个系数。总之
\begin{equation}
    \label{Eq:I:42:2}
    N_c=nv
\end{equation}
是到达每单位面积并凝聚的分子数。

同时，液体的分子不停地跳动，时而它们中的某个分子被“踢”出液面。现在我们来估计它们被踢出液面有多快。这里的概念是，在平衡时每秒钟踢出的分子数与到达的分子数相等。

踢出的分子有多少？为了能跑出液面，必须要求某个特殊的分子偶然获得超过其邻近分子的额外能量——一个可观的额外能量，因为它受到液体中其他分子的很强的吸引力。通常由于这种很强的吸引力分子不会离开液体，但是在碰撞中有时它们中的某个分子会偶然得到一份额外的能量。若$W\gg kT$ ，那么要得到在我们的情况下所需的额外能量$W$的机会是非常小的。事实上，$e^{-W/kT}$ 是一个分子得到大于这么多的能量的概率。这是分子动理论的一条一般性原理：为了取得超过平均值的额外能量$W$，概率是$e$的负指数幂，指数是我们所要获得的能量除以$kT$。现在假设有一些分子获得了这个能量。我们来估计每秒钟有多少分子离开液面。当然，仅仅使分子具有必要的能量并不意味着它实际上必然蒸发，因为它可能浸没在液体内很深的地方，或者，即使它接近表面，但可能沿着不适当的方向运动。每秒钟离开单位面积的分子数应当是：靠近表面的单位面积的分子数除以使得一个分子离开所需的时间，然后再乘上它们具有足够能量跑出液面的概率$e^{-W/kT}$。

我们假设液面上的每个分子占有一定的截面$A$ ，则液面上单位面积内的分子数为$1/A$。现在一个分子跑出液面需要多长的时间？如果分子具有一定的平均速率$v$，以及足够的能量可以移动一个分子的直径$D$那样的路程（即第一层的厚度），则它经过这一厚度的时间就是离开液面所需的时间。这段时间即为$D/v$。于是蒸发的分子数近似地可写为
\begin{equation}
    \label{Eq:I:42:3}
    N_e=(1/A)(v/D)e^{-W/kT}
\end{equation}

式中每个分子的面积乘以这层分子的厚度近似地等于单个分子所占的体积$V_a$，因此为了达到平衡，必须使$N_c=N_e$，或者
\begin{equation}
    \label{Eq:I:42:4}
    nv=(v/V_a)e^{-W/kT}
\end{equation}

我们可以消去$v$，因为它们相等，虽然一项是蒸气中分子的速度，另一项是正在蒸发的分子的速度，但它们是相同的，因为我们知道它们在一个方向上的平均动能是$kT/2$。但是有人会反驳说：“不！不！这些是运动得特别快的分子，是获取了额外能量的分子。”不过，实际上并非如此，因为开始把它们从液体中拉出来时，为了抵抗势能，它们必须消耗 这些额外能量。所以，当它们跑到液面时，已经减慢到速度$v$！它与我们讨论过的大气中分子速度分布的情况相同：在底部，分子具有一定的能量分布，而到达顶部的那些分子具有同样 的能量分布，因为速度慢的分子根本到达不了上面，而速度快的分子的运动则减慢下来。正在蒸发的分子和液体内的分子的能量分布相同——这是一个颇为引人注目的事实。无论如何，因为还有种种其他的不精确性，诸如存在着蒸气分子从液面弹回而不是跑入液体的概率，等等，所以过于细致地讨论我们的公式是无益的。这样我们对于蒸发率与凝结率有了一个粗略的概念，当然，我们看到蒸气密度$n$变化的方式还是像以前那样，但是现在我们比较具体地理解了它，而不是只把它看成一个任意的公式。

这个比较深入的理解使我们得以分析某些事情。例如，假定我们抽出蒸气的速率快到跟蒸气形成的速率一样（倘若有一个很好的泵，而液体的蒸发又很慢），那么，如果将液体温度保持为$T$，蒸发会有多快？假设我们已从实验上测量了处于平衡状态时的蒸气密度，因而知道在给定的温度下，与液体处于平衡状态时单位体积的分子数。现在我们希望知道液体蒸发得有多快 。虽然迄今为止我们只对它蒸发的部分作过粗略的分析，但除去有一个未知的反射系数因子外，到达 液面的蒸气分子数也可以大体知道。所以我们可以利用平衡时离开液面的分子数与到达液面的分子数相等的事实。确实，当不断抽去蒸气时，分子只会跑出液面。但是如果蒸气被留下来，它将达到平衡状态下的密度，这时跑回液体的分子数等于蒸发的分子数。因此，不难看出每秒离开液面的分子数等于某个未知的反射系数$R$乘以假设蒸气仍然存在时每秒返回液面的分子数，因为这就是在平衡时应该补偿掉的数目，故
\begin{equation}
    \label{Eq:I:42:5}
    N_e=(vR/V_a)e^{-W/kT}
\end{equation}
当然，从蒸气中碰撞到液体上的分子数是不难算出的，因为我们毋须像在考虑分子如何离开液面那样，要对有关的力知道得那么多，用其他方式作论证要容易得多。

\section{热离子发射}

我们举另外一个很实际的例子，它和液体的蒸发非常类似，以至于不值得另行分析。这在本质上是同样的问题。在电子管中，有一个电子发射源，即一根加热的钨丝以及一块吸引电子的带正电的板极。任何电子离开钨丝表面后就直接向板极飞去。这是一个理想的“泵”，它不停地“抽掉”电子。现在的问题是：每秒有多少电子从钨丝发射出来？而这个电子数目又怎样随温度而变化？这个问题的答案与式（\ref{Eq:I:42:5}）相同，因为在一块金属中，电子受到金属的离子或原子的吸引。粗略地讲，它们被金属吸引。为了从金属中拉出一个电子，必须提供一定的能量，或拉出它来的功。这个功的大小视金属的种类而异。事实上，它甚至随某种给定金属表面的性质而变化，但是总的功约为几个电子伏，附带提一下，它是化学反应中所涉及的典型的能量。只要回忆一下化学电池中由化学反应产生的类似闪光灯电池上的电压约为1V左右就可以想起这一事实了。

我们怎么能够找出每秒钟有多少电子跑出来呢？分析电子跑出的效应是十分困难的，用其他方法来分析比较容易。所以，开始时我们可以设想并没有将电子抽走，电子就像气体一样可以跑回到金属内。因此在平衡时应当有一定的电子密度，当然，它可以准确地由与式（\ref{Eq:I:42:1}）同样的公式给出。在这里，粗略地讲，$V_a$是金属中每个电子的体积，而$W$等于$q_e \phi$ ，其中$\pi$称为功函数 ，或者说是把一个电子拉出金属表面所需的电势。这就告诉我们，应当有多少电子存在于金属周围的空间中并撞到金属上，以平衡从金属跑出来的电子。于是，很容易算出如果抽走所有电子的话，会有多少电子从金属跑出，因为跑出的电子数恰好等于具有上述密度的电子“蒸气”跑回金属的电子数。换句话说，答案是：单位面积流入的电流等于每个电子的电荷乘以每秒到达单位面积的电子数，后者正如我们已多次见到的那样，就是单位体积的电子数乘以速度
\begin{equation}
    \label{Eq:I:42:6}
    I=q_env=(q_ev/V_a)e^{-q_e\phi/kT}
\end{equation}
而1 eV相当于温度为11 600 K的$kT$值。电子管灯丝比方说可以在温度为1 100 K下工作，这样指数因子约为$e^{-10}$，当我们改变一点温度时，指数因子变化很大。这样，我们再次看到公式的主要的特征是 。实际上，前面的因子是完全不正确的——结果表明金属中的电子的行为不能从经典理论得到正确描写，而要用量子力学才行，但是这只使前面的因子改变一点点。事实上，还没有一个人能够很好地直接求得它，虽则为了计算它许多人应用了高等量子力学理论。主要的问题是，$W$是否略微随$T$而变化？如果是的话，人们就无法把$W$随温度的缓慢变化和式前系数上的差别区别开来。比方说，如果$W$随温度呈线性变化，那么$W=W_0+\alpha kT$，这样就有
\begin{equation*}
    e^{-W/kT} = e^{-(W_0+\alpha kT)/kT} = e^{-\alpha}\cdot e^{-W_0/kT}
\end{equation*}
于是，$W$和温度的线性依赖关系等价于“常数”的改变。试图精确地得到式前的系数确实是非常困难的，并且通常是徒劳的。

\section{热电离}

我们现在继续讨论另外一个应用同样概念的例子——总是同样的概念。这个例子与电离有关。假设在一种气体中有大量原子，它们处在中性状态，但气体很热，原子可能被离子化。我们想知道，如果温度一定，并且每单位体积的原子密度一定，在这种情况下将有多少离子。我们仍旧考虑一个容器，其中共有$N$个带有电子的原子（如果一个电子离开一个原子，那个原子就叫做离子，而如果原子呈中性，我们就简单地叫它为原子）。假设在任何确定的时刻，每单位体积中的中性原子数是$n_a$，离子数是$n_i$，而电子数是$n_e$。问题是：这三个数之间的关系如何？

首先，在这些数之间存在两个条件或约束。例如，在改变不同的条件，像温度或其他时，$n_a+n+n_i$将保持不变，因为这正是容器中的原子核数$N$。如果使单位体积的原子核数保持固定，而比方说改变温度，那么在电离过程中，某些原子将变为离子，但原子加离子的总数不会变化，也就是说
\begin{equation*}
    n_a+n_i=N
\end{equation*}
另一个条件是，如果整个气体呈中性（如果我们忽略二次或三次电离），这就意味着在所有时间内离子的数目等于电子的数目，即
\begin{equation*}
    n_i=n_e
\end{equation*}
这就是两个补充方程，它们只是简单地表示电荷守恒与原子守恒。

这些方程是正确的，我们在考虑实际问题时，最后要用到它们。但我们要想得到这些量之间的另一个关系式。为此可以这样做，我们再用类似的概念：电子离开原子需要一定的能量，称为电离能。同时为了使所有公式看起来一样，我们将它写为$W$。这样，$W$就是从原子中拉出一个电子以产生一个离子所需要的能量。我们仍认为，在“蒸气”中每单位体积的自由电子数等于原子中单位体积的束缚电子数乘以$e$的负指数幂，其指数为束缚态与自由态之间的能量差除以$kT$。这又是一个基本方程。我们怎么把它写出来呢？每单位体积的自由电子数当然是$n_e$，因为这就是$n_e$的定义。那么，每单位体积中束缚在原子中的电子数是多少？可以放进电子的位置的总数看来是$n_a+n_i$，我们假设电子被束缚时，每个电子都束缚在一定的体积$V_a$之内。所以总的可束缚电子的体积将是$(n_a+n_i)V_a$，这样我们可以把公式写为
\begin{equation*}
    n_e=\frac{n_a}{(n_a+n_i)V_a}e^{-W/kT}
\end{equation*}
但这个式子在一个基本方面是错误的，即当一个电子已经在一个原子内以后，另一个电子再也不可能进入这个体积之内！换句话说，对于一个正在考虑究竟是要跑到“蒸气”中去，还是要凝结在某一位置上的电子来说，实际上并非所有可能位置的全部体积都是可供选择的。因为在这个问题中还有另外的特点：当已有其他电子处在那里时，这个电子就不能跑进去——它受到排斥力。由于这个原因，我们应当计算的只是电子可能位于那里或者不位于那里的那部分体积。也就是说，已为电子占据的体积不计在总的可供选择的体积之内，而允许占有的只是离子 的体积，在那里有空位可让电子占据。所以，在这种情况中，我们发现可将公式写得更准确一些，就是
\begin{equation}
    \label{Eq:I:42:7}
    \frac{n_en_i}{n_a} = \frac{1}{V_a}e^{-W/kT}
\end{equation}
这个公式叫作萨哈（Saha）方程 。现在我们来看看能否从所发生的动力学讨论中定性地了解为什么类似于这样的公式是正确的。

首先，每当一个电子跑到离子内，它们就组合成一个原子。同样，每当一个原子受到碰撞后，就会分裂为一个离子和一个电子。这两种比率应该相等。电子与离子彼此发现对方有多快呢？如果每单位体积的电子数增加，这种比率一定增加。如果每单位体积的离子数增加，比率也会增加。这就是说，产生复合的总比率一定正比于电子数和离子数的乘积。而由于碰撞引起的总电离率必定线性地依赖于电离的原子有多少，于是，当$n_en_i$的乘积与原子数$n_a$之间有某种关系时，两种比率将得到平衡。这种关系正好能用式（\ref{Eq:I:42:7}）这个特殊公式（其中$W$是电离能）表述出来，这一事实当然包含的内容要更多一些。但我们不难理解公式必然包含像$n_e n_i /n_a$这样的电子浓度、离子浓度与原子浓度之间的组合以得到一个与各个密度无关而只取决于温度、原子截面和其他的常数因子的常数。

我们可能还注意到，由于方程中包括每单位体积的粒子数，如果在原子加离子的总数$N$固定，即原子核的总数固定情况下做两个实验，但用两个体积不同的容器，则在大容器中，所有的$n$都较小。但由于比值$n_e n_i /n_a$保持不变，因此，在较大的容器中电子与离子的总数必然较大。为了看出这一点，假定在体积为$V$的容器内有$N$个原子核，其中被电离的百分比是$f$。则$n_e =fN /V =n_i$，而$n_a =(1-f )N /V$。这样我们的方程变为
\begin{equation}
    \label{Eq:I:42:8}
    \frac{f^2}{1-f}\frac{N}{V} = \frac{e^{-W/kT}}{V_a}
\end{equation}
换句话说，如果我们使原子密度越来越小，或使容器体积越来越大，电子与离子的百分比$f$必定增加。由于存在这种只是因密度下降即“膨胀”引起的电离现象，使得我们即使还不能从适当的能量观点来理解它，但也有理由相信，在极低的密度下，比如在恒星间的冷的空间中，也可能存在着离子。虽然为了产生离子需要许多$kT$的能量，但离子还是存在。

为什么当周围空间这么大时也会有离子存在，而当我们增加密度时，离子却趋于消失？答案是：考虑一个原子。偶尔，会有光，会有另一个原子或离子，总之会有保持热平衡的某个东西去撞击它。这种情况是极其罕见的，因为要使一个电子飞去而留下一个离子，需要大量的额外能量。如果空间极大，那么电子到处荡来荡去，经过年复一年也许还没有接近任何东西。但极其偶然地，它回到一个离子处，并和离子复合而组成一个原子。所以电子从原子中跑出来的比率非常低。但是因为空间体积极大，一个已经跑出来的电子要经过很长时间才能找到另一个和它复合的离子，因而复合概率非常小。因此，尽管需要很大的额外能量，仍然会有一定数量的电子存在。

\section{化学动力学}

在化学反应中也可以见到我们刚才称为“电离”的同样情况。例如，假定两个物体A与B组成化合物AB，稍加思考后就能看出，AB就是我们所谓的“原子”，B可以叫作“电子”，A可以叫作“离子”。经过这样代替以后，平衡方程式的形式就完全相同
\begin{equation}
    \label{Eq:I:42:9}
    \frac{n_An_B}{n_{AB}}=Ce^{-W/kT}
\end{equation}
当然，这个公式不是严格的，因为“常数”$C$与容许A和B组合的体积等有关。但是，利用热力学的论据，人们可以确定指数因子上的$W$的含义，结果表明它与反应中所需的能量非常接近。

假定我们试图把这个公式理解为碰撞的结果，就像我们在理解蒸发公式时那样，看看单位时间有多少电子跑出来，又有多少电子跑回去。假设A和B在偶尔进行的一次碰撞中形成一个化合物AB。假设AB是一个复杂分子，它四面八方作无规则运动，不时受到其他分子的碰撞，使它逐渐得到足够的能量以重新分裂为A和B。

实际表明，在化学反应中，如果原子靠近时带有的能量太小，即使在$A+B\longrightarrow AB$的反应中可能释放能量，但是，A与B可能彼此接触这一事实不一定会使反应开始进行。事实上，通常要求碰撞得很厉害才能使反应得以发生，A与B之间的一次“软”碰撞不可能做到这一点，即使在反应过程中可能释放能量。所以我们假定在化学反应上这样的情况是非常普遍的：为了使A与B形成AB，它们不能只靠彼此碰撞，而必须在碰撞时具有足够的能量。这个能量称为激活能——亦即使反应“激活”所需的能量。我们称$A^*$为激活能，为了使反应实际上得以进行，这是碰撞中所要求的额外能量。A和B产生AB的速率$R_f$应当包括A的原子数乘B的原子数，乘以单个原子碰撞在一定截面$\sigma_{AB}$上的速率，再乘以因子$e^{-A^*/kT}$，$e^{-A^*/kT}$是它们具有足够能量的概率
\begin{equation}
    \label{Eq:I:42:10}
    R_f=n_An_Bv\sigma_{AB}e^{-A^*/kT}
\end{equation}
现在我们要求出相反的速率$R_r$。AB也有一定的飞散开的机会。为了能够飞散开，不仅必须有一个使A、B完全分开的能量，而且正像使A和B难以化合在一起那样，必须存在一类A和B必须爬过它才能再次分开的能量“小丘”，它们必须不但具有正好可使彼此分开的足够能量，而且必须有一定的额外能量。这就像要翻过一座小山进入到一个深
\begin{wrapfigure}{l}{0.4\textwidth}
    \centering
    \includegraphics[width=0.38\textwidth]{Chapter42/化学反应A+B——AB的能量关系}
    \caption{\footnotesize 化学反应$A+B\longrightarrow AB$的能量关系}
    \label{figure:化学反应A+B——AB的能量关系}
\end{wrapfigure}
谷一样，人们得爬过小山进入深谷，在返回时又得爬出深谷，然后再爬过小山（图\ref{figure:化学反应A+B——AB的能量关系}）。于是AB变成A和B的速率将正比于已存在的数目$n_{AB}$，再乘以$e^{-(W+A^*)/kT}$
\begin{equation}
    \label{Eq:I:42:11}
    R_r=C'n_{AB}e^{-(W+A^*)/kT}
\end{equation}
这里的$C'$与原子的体积和碰撞率有关，我们可以像在蒸发的情况下所做的那样，用面积、时间和厚度来求出它，但我们在这里将不这样做。我们感兴趣的主要事情是在这两种速率相等时，它们的比值等于1。这告诉我们
\begin{equation*}
    \frac{n_An_B}{n_{AB}}=Ce^{-W/kT}
\end{equation*}
像以前一样，这里的$C$包括截面、速度以及其他与$n$无关的因子。

有趣的是，反应率也随$exp[-(\text{常数})/(kT)]$而变化，虽然这里的常数与支配浓度的常数并不相同，激活能$A^*$ 与能量$W$也截然不同。$W$支配在平衡状态下A、B与AB的比例 ，但是，如果我们想知道A+B变为AB有多快，这就不是一个平衡的问题了，这里有一种不同的能量——激活能 ，通过指数因子来支配反应率。

此外，$A^*$并不是一个像$W$那样的基本常数。假如在器壁的表面上——或在某些其他地方——A和B得以按某种方式暂时粘附在那里，从而能使它们更易化合。换句话说，我们可以找到一条穿过小山的“隧道”，或者一座较低的小山。由能量守恒定律知道，当一切都完成后，我们仍由A和B组成AB，因而能量差$W$与反应发生的方式完全无关，但是激活能$A^*$则与反应发生的方式极有关系 。这就是为什么化学反应率对外部条件非常敏感的原因。如果化学反应依赖于表面的性质，我们可以通过放入一个不同类型的表面以改变反应率，也可以使反应率在一个“不同的圆筒”中进行，这样它就有不同的反应率。或者，如果我们放入第三种物质，它可以大大改变反应率。有的东西仅仅通过使$A^*$稍加变化就能使反应率产生极大的变化——我们把它叫做催化剂 。在一定的温度下，由于$A^*$过大，可能会使一个反应实际上根本无法进行，但是，当加入这种特殊物质——催化剂后，因为降低了$A^*$，反应就进行得非常快。

附带说一下，在A加B产生AB这样的反应中，存在着一些麻烦，因为当我们企图将两种物质放在一起以组成一种更稳定的物质时，无法使能量和动量两者都守恒。因此，我们至少还需要第三种物质C，这样实际反应就远为复杂了。正方向的反应速率包括乘积$n_A n_B n_C$，看来，我们的公式是错误的，但并非如此！当我们考察AB沿反方向的反应率时，我们发现它也需要和C发生碰撞，这样，在反方向的反应率中也存在着$n_{AB}n_C$，而在浓度平衡公式中$n_C$已被消去。无论反应机制如何，我们开始写下的平衡定律式（\ref{Eq:I:42:9}）绝对保证正确！

\section{爱因斯坦辐射律}

我们现在转到一个类似的有趣问题，它与黑体辐射定律有关。在上一章中，我们用普朗克的方法，考虑了振子的辐射问题，得出了空腔内的辐射分布定律。振子必定具有某个平均能量，既然它在振动，它将辐射，而且不断地往空腔辐射能量，直至它积蓄了足够的辐射以使吸收和发射保持平衡。这样我们得到频率为$\omega$的辐射强度由下列公式给出
\begin{equation}
    \label{Eq:I:42:12}
    I(\omega)d\omega = \frac{\hbar\omega^3d\omega}{\pi^2c^2(e^{\hbar\omega/kT}-1)}
\end{equation}
这个结果包括了进行辐射的振子具有有限的、间隔相等的能级的假设。我们没有说过光必定是光子或任何某种与之类似的东西。当一个原子从某一能级跑到另一个能级时，能量必然以光的形式按一个能量单位$\hbar\omega$释放出来。至于为什么会这样我们还没有讨论过。普朗克的原始思想认为物质而不是光是量子化的：物质振子不能取任何能量，它必须一份一份地取得能量。此外，推导中的麻烦还在于它是部分经典的。我们按经典物理计算了一个振子的辐射率，然后又回过来说：“不，这个振子有一系列能级。”所以，为了求得正确的和完全量子力学化的结果，物理学经历了非常缓慢的发展，终于在1927年建立了量子力学。但在此期间，爱因斯坦作了努力，他改变了普朗克的只有物质振子才可量子化的观点，而认为实际上光是一些光子，并且在某些方面可以看作是具有能量$\hbar\omega$的粒子。此外，玻尔指出，任何 原子系统都有能级，但并不需要像普朗克振子那样均匀地隔开。所以，从更完全的量子力学观点来重新推导，或者至少重新讨论辐射定律就成为必要了。

爱因斯坦假设普朗克的最后公式是正确的，他利用这个公式得到了某些前所未知的有关辐射和物质相互作用的崭新的知识。他的讨论如下：考虑一个原子的许多能级中的
\begin{wrapfigure}{l}{0.38\textwidth}
    \centering
    \includegraphics[width=0.35\textwidth]{Chapter42/原子内两个能级间的跃迁}
    \caption{原子内两个能级间的跃迁}
    \label{figure:原子内两个能级间的跃迁}
\end{wrapfigure}
任何两个能级，比方说第$m$个与第$n$个能级（图\ref{figure:原子内两个能级间的跃迁}）。爱因斯坦假设，当这样一个原子受到适当频率的光照射时，它能够吸收光线中的光子，使它从$n$态跃迁到$m$态，而每秒钟发生这种情况的概率当然与这两个能级有关，但是它也正比 于照射在原子上的光线有多强 。我们把比例常数称为$B_{nm}$ ，以便提醒我们这不是一个自然界的普适常数，而只是与特定的两个能级有关：有的能级易于激发，有的能级则难于激发。现在，从$m$到$n$的发射率公式是什么？爱因斯坦认为必须把它分为两部分。首先，即使没有光存在，也应当存在原子从激发态落到较低能态而发射出一个光子的机会，我们称这一现象为自发发射。这与一个具有一定能量的振子即使在经典物理中也不会保持能量，而将通过辐射丧失能量的思想相似。于是，经典系统自发辐射的类比是：如果原子处在一种激发态，那么它就存在着一个确定的从$m$下降到$n$的概率$A_{mn}$，$A_{mn}$仍是依赖于这两个能级，但与原子是否受到光的照射无关。但是爱因斯坦再前进了一步，通过与经典理论的比较以及其他的论据得出，发射也受光的存在的影响——当恰当频率的光照射到原子上时，它会增加一个光子的发射率，而这个发射率正比于光的强度，比例常数是$B_{mn}$。如果后来我们推导出这个系数是零的话，就能发现爱因斯坦错了；当然，我们将发现他是对的。

这样，爱因斯坦假定存在着三种过程：一种是正比于光强度的吸收，一种是正比于光强度的发射，它称为感生发射或者有时称为受激发射，还有一种是与光无关的自发发射。

现在我们假设在温度为$T$的平衡条件下，在$n$态有$N_n$个原子，在$m$态有$N_m$个原子。那么由$n$态跑到$m$态的原子总数就是在$n$态的原子总数乘以每秒钟一个原子从$n$态上升到$m$态的比率。所以我们就得到一个每秒钟从$n$上升到$m$的原子数的公式
\begin{equation}
    \label{Eq:I:42:13}
    R_{n\rightarrow m}=N_nB_{nm}I(\omega)
\end{equation}
从$m$到$n$的数目可以用同样方式表示，即在$m$态的原子数$N_m$乘以每秒钟一个原子掉到$n$态的概率。这一次我们的表示式是
\begin{equation}
    \label{Eq:I:42:14}
    R_{n\rightarrow m}=N_m[A_{mn}+B_{mn}I(\omega)]
\end{equation}

现在我们假设，在热平衡时跑上去的原子数必定等于落下来的原子数。至少，这是使每个能级上的原子数确实保持不变的一种方式\footnote{这并非唯一可能的使不同能级上的原子数保持不变的安排方式，但这是实际发生的情况。在热平衡状态下，每一个过程必定被它准确的相反过程所平衡，这一原理称为细致平衡原理。}。所以我们将假定这两种比率在平衡时相等。此外，我们还有另一种情况：我们知道$N_m$与$N_n$相比有多大——两者的比值是$e^{-(E_m-E_n)/kT}$。现在爱因斯坦假设：从$n$到$m$的跃迁中唯一有效的光线只是具有相应于能级差的频率的光，所以在我们的所有公式中就有$E_m-E_n=\hbar\omega$，于是
\begin{equation}
    \label{Eq:I:42:15}
    N_m=N_ne^{-\hbar\omega/kT}
\end{equation}
这样，如果使两种比率相等
\begin{equation*}
    N_n B_{nm} I(\omega)=N_m [A_{mn} +B_{mn}I(\omega)]
\end{equation*}
除以$N_m$，得到
\begin{equation}
    \label{Eq:I:42:16}
    B_{nm} I(\omega)e^{\hbar\omega/kT}=A_{mn} +B_{mn}I(\omega)
\end{equation}
我们可以由这个方程算出$I(\omega)$，它就是
\begin{equation}
    \label{Eq:I:42:17}
    I(\omega)=\frac{A_{mn}}{B_{nm}e^{\hbar\omega/kT}-B_{mn}}
\end{equation}
但普朗克已经告诉我们这个公式必须是式（\ref{Eq:I:42:12}）。由此我们可以推出一些结论：首先，$B_{nm}$必定等于$B_{mn}$，否则我们就不能得出$(e^{\hbar\omega/kT}-1)$。所以，爱因斯坦发现了某种他还不知道如何计算的东西，即感生发射概率必定等于吸收概率 。这是十分有趣的。此外，为了使式（\ref{Eq:I:42:17}）与式（\ref{Eq:I:42:12}）一致，$A_{mn} /B_{mn}$ 必须为 
\begin{equation}
    \label{Eq:I:42:18}
    \frac{A_{mn}}{B_{mn}}=\frac{\hbar\omega^3}{\pi^2c^2}
\end{equation}
所以，比如说如果已知一定能级的吸收率，就可以推导出自发发射率和感生发射率，或者它们的任何组合。

这就是迄今为止爱因斯坦或任何别的人利用这样的论证所能得到的结果。当然，要实际计算绝对自发发射率或对任何特殊原子跃迁的其他比率，都需要所谓量子电动力学的有关原子内部机制的知识，这些知识直到11年后才被人们所发现。而爱因斯坦的这个工作是在1916年完成的。

今天感生发射的可能性已找到了颇有意义的应用。如果有光存在，它将倾向于激发向下的跃迁。如果某些原子处于较高的状态，跃迁就会使可资应用的光能增加$\hbar\omega$。现在，通过某种非热学的方法，可以使气体中$m$态的原子数远大于$n$态的原子数。这种情况远离平衡态，因此不能用平衡态的公式$e^{-\hbar\omega/kT}$。我们甚至可以安排得使较高状态的原子数非常之大，而较低状态的原子数实际为零。于是具有相应于能量差$E_m-E_n$的频率的光不会强烈地被吸收，因为在$n$态中没有那么多原子去吸收它。另一方面，当这样的光存
\begin{wrapfigure}{l}{0.4\textwidth}
    \centering
    \includegraphics[width=0.38\textwidth]{Chapter42/如用蓝光将原子激发到较高的状态h，它可以发射一个光子使原子处于m态，m态的原子足够多时就开始产生激光}
    \caption{\footnotesize 如用蓝光将原子激发到较高的状态$h$，它可以发射一个光子使原子处于$m$态，$m$态的原子足够多时就开始产生激光}
    \label{figure:如用蓝光将原子激发到较高的状态h，它可以发射一个光子使原子处于m态，m态的原子足够多时就开始产生激光}
\end{wrapfigure}
在时，它将诱发从高能态的发射！所以，如果在高能态有大量原子，就会产生一种连锁反应，也就是说一旦原子开始发射，越来越多的其他原子将引起发射，最后大量的原子都一起倾泻到低能级上去。这就是所谓的激光 ，或者在远红外的情况下，就称为微波激射 （图\ref{figure:如用蓝光将原子激发到较高的状态h，它可以发射一个光子使原子处于m态，m态的原子足够多时就开始产生激光}）。

为了得到处在$m$态的原子，可以应用各种技巧。如果我们用高频的强光束照射，原子就能到达更高的能级。它们能够从这些高能级上发射出各种光子而落下来，直至全部到达$m$态。如果它们趋向于停留在$m$态而不发射，这种状态称为亚稳态。以后，通过感生发射，它们一起倾泻到$n$态。一个比较技术性的问题是，假如我们将这个系统放在一个普通的容器内，那么与感应效应相比较，它们将沿许多方向自发辐射，以致又会引起麻烦。但是，如果在容器的每一边放上一面接近于理想的镜子，以致被发射的光得到更多的机会一而再、再而三地感生更多的发射，那么我们就能够加强感应效果，提高它的效率。虽然镜子几乎能100\%地反射，但它仍有少许透射，因而有一点点光跑出去。当然，最后从能量守恒定律看来，所有的光都沿一个统一的笔直的方向跑出去，而形成强光束。今天，利用激光器可能做到这一点。

